%% demo.tex -- Showcase deck for the tuftedark beamer theme.
%% Compile from this directory:  latexmk -pdf demo.tex   (or pdflatex x2)
%% Exercises every layout, plus one real equation and one real dark-native
%% figure adapted from the Unsupervised Learning lecture notes.

\documentclass[aspectratio=169]{beamer}
\usetheme{tuftedark}

\title[Unsupervised Deep Learning]{Unsupervised Deep Learning}
\subtitle{A Tufte-dark slide template}
\author[Schutera]{Prof.\ Dr.-Ing.\ Mark Schutera}
\institute{DHBW Ravensburg}
\date{\today}

% auto section-divider slides
\AtBeginSection[]{\begin{frame}\sectionpage\end{frame}}

\begin{document}

\begin{frame}[plain]
  \titlepage
\end{frame}

% --- Agenda: the "lecture in one slide" (why it matters + what you'll learn) ---
\begin{frame}{Why this lecture}
  \begin{withmargin}
    \begin{tdmain}
      Most data arrives without labels. If structure only existed where someone
      drew it, we would be blind to most of what is there.
      \par\bigskip
      \centering
      \begin{tikzpicture}[node distance=8mm]
        \node[tdbox] (a) {unlabelled\\data};
        \node[tdbox, right=of a] (b) {hidden\\structure};
        \node[tdboxhi, right=of b] (c) {decisions};
        \draw[tdarr] (a) -- (b);
        \draw[tdarr] (b) -- (c);
      \end{tikzpicture}
    \end{tdmain}
    \begin{tdmargin}
      \textbf{After this lecture}
      \par\smallskip
      you can find groups without labels, judge how many there are, and read
      what they mean.
    \end{tdmargin}
  \end{withmargin}
\end{frame}

% ======================================================================
\section{Foundations}
% ======================================================================

\begin{frame}{What clustering buys us}
  Given unlabelled data $\{x_1,\dots,x_n\}$, we want structure without
  supervision: a partition into $K$ groups that are internally coherent.
  There is no ground truth, so the objective \alert{is} the supervision,
  and our choice of distance and of $K$ encodes the prior over what
  ``coherent'' should mean.
  \par\medskip
  Two failure modes follow directly:
  \begin{itemize}
    \item Euclidean distance silently assumes isotropic, comparable scales.
    \item The wrong $K$ hides real structure, or invents structure that is not there.
  \end{itemize}
\end{frame}

\begin{frame}{Margin notes}
  \begin{withmargin}
    \begin{tdmain}
      Lloyd's algorithm alternates assignment and update until the partition
      stops changing. Each step cannot increase $\mathcal{J}$, so it
      converges, though only to a local minimum.
      \par\medskip
      The centroid update is just the per-cluster mean: the choice that
      minimises that cluster's contribution to the objective.
    \end{tdmain}
    \begin{tdmargin}
      Lloyd, 1957; published 1982.
      \par\smallskip
      The margin holds the aside, the citation, the small caveat: present,
      but out of the main line of thought.
    \end{tdmargin}
  \end{withmargin}
\end{frame}

\begin{frame}{The $k$-means objective}
  We minimise within-cluster squared distance to the centroids:
  \[
    \mathcal{J}(C,\mu)
      \;=\; \sum_{k=1}^{K}\;\sum_{x \in C_k}
            \bigl\lVert\, x - \textcolor{DHBWred}{\mu_k}\,\bigr\rVert^{2},
    \qquad
    \textcolor{DHBWred}{\mu_k} = \frac{1}{|C_k|}\sum_{x\in C_k} x .
  \]
  Ward linkage extends the same variance idea to merges:
  \[
    D_{\mathrm{Ward}}(C_i,C_j)
      \;=\; \frac{|C_i|\,|C_j|}{|C_i|+|C_j|}\,
            \bigl\lVert \mu_i - \mu_j \bigr\rVert^{2}.
  \]
  The accent marks the centroid: the one quantity every step recomputes.
\end{frame}

% --- Equation on the 2/3 + 1/3 grid: artefact left, reading right ---
\begin{frame}{Reading the objective}
  \begin{withmargin}
    \begin{tdmain}
      \vspace*{1.2em}
      \[
        \mathcal{J}(C,\mu)
          \;=\; \sum_{k=1}^{K}\;\sum_{x \in C_k}
                \bigl\lVert\, x - \textcolor{DHBWred}{\mu_k}\,\bigr\rVert^{2}
      \]
    \end{tdmain}
    \begin{tdmargin}
      The summed squared distance from every point to \textcolor{DHBWred}{its own
      centroid}. Driving it down is precisely what we mean by a coherent cluster.
    \end{tdmargin}
  \end{withmargin}
\end{frame}

% --- Derivation slide: multiple steps, smaller font, no side column ---
\begin{frame}{Derivation: one Ward merge}
  \small
  \begin{align*}
    D_{\mathrm{Ward}}\!\left(\{x_1,x_2\},\, x_3\right)
      &= \frac{n_{12}\, n_3}{n_{12}+n_3}\,\bigl\lVert \mu_{12} - x_3 \bigr\rVert^{2} \\[2pt]
      &= \frac{2 \cdot 1}{2 + 1}\,\bigl\lVert (2,1) - (8,7) \bigr\rVert^{2} \\[2pt]
      &= \tfrac{2}{3}\,\bigl\lVert (-6,-6) \bigr\rVert^{2} \\[2pt]
      &= \tfrac{2}{3}\cdot 72 \;=\; 48
  \end{align*}
\end{frame}

\begin{frame}{Definition block}
  \begin{tddef}{Centroid}
    The centroid $\mu_k$ of cluster $C_k$ is the mean of its members.
    It is the point that minimises the summed squared distance within the
    cluster, which is exactly the per-cluster term of $\mathcal{J}$.
  \end{tddef}
  \vskip1ex
  \begin{tdblock}
    A plain statement block: same rounded, thick-bordered vocabulary as the
    print TikZ boxes, no title bar.
  \end{tdblock}
\end{frame}

\begin{frame}{Blocks}
  \begin{block}{Standard block}
    The native beamer block, rounded with a grey title bar (KIT structure,
    DHBW colours). For definitions, statements, summaries.
  \end{block}
  \begin{exampleblock}{Example block}
    A second tone for worked examples and asides.
  \end{exampleblock}
  \begin{alertblock}{Alert block}
    Red, reserved: a warning or the one thing not to miss.
  \end{alertblock}
\end{frame}

\begin{frame}{Two columns}
  \begin{columns}[T]
    \begin{column}{0.5\textwidth}
      \textbf{Lloyd's algorithm}
      \begin{itemize}
        \item Assign each $x$ to nearest $\mu_k$.
        \item Recompute each $\mu_k$.
        \item Repeat until assignments settle.
      \end{itemize}
    \end{column}
    \begin{column}{0.5\textwidth}
      \textbf{Caveats}
      \begin{itemize}
        \item Converges to a local minimum.
        \item Sensitive to initialisation.
        \item Assumes roughly spherical clusters.
      \end{itemize}
    \end{column}
  \end{columns}
\end{frame}

\begin{frame}{Figure: the elbow (dark-native tufteplot)}
  \centering
  \begin{tikzpicture}
    \begin{axis}[tufteplotdark, width=0.72\textwidth, height=0.40\textwidth,
        xlabel={number of clusters $K$}, ylabel={inertia $\mathcal{J}$},
        xtick={1,2,3,4,5,6}, ytick={0,40,80,120},
        ymin=0, ymax=130, xmin=0.6, xmax=6.4]
      \addplot[mark=*, mark size=1.2pt, thick, tdfg]
        coordinates {(1,120)(2,72)(3,40)(4,30)(5,24)(6,20)};
      \addplot[mark=*, mark size=2.2pt, only marks, DHBWred]
        coordinates {(3,40)};
    \end{axis}
  \end{tikzpicture}
  \par\smallskip
  {\small\color{tddim}The marginal gain flattens past $K=3$ (accent): the elbow.}
\end{frame}

\begin{frame}{Figure: a pipeline (dark-native TikZ)}
  \centering
  \begin{tikzpicture}[node distance=8mm]
    \node[tdbox] (x)   {input $x$};
    \node[tdbox, right=of x] (enc) {encoder};
    \node[tdboxhi, right=of enc] (z) {latent $z$};
    \node[tdbox, right=of z] (dec) {decoder};
    \node[tdbox, right=of dec] (xh) {$\hat{x}$};
    \draw[tdarr] (x) -- (enc);
    \draw[tdarr] (enc) -- (z);
    \draw[tdarr] (z) -- (dec);
    \draw[tdarr] (dec) -- (xh);
    \draw[tdarrdim] (xh.south) to[out=-90,in=-90] node[below,td,font=\scriptsize]{reconstruction loss} (x.south);
  \end{tikzpicture}
\end{frame}

% ======================================================================
\section{Results}
% ======================================================================

\begin{frame}{Comparison table}
  \centering
  {\arrayrulecolor{tddim}%
  \begin{tabular}{lcc}
    \toprule
    method & needs $K$ & cluster shape \\
    \midrule
    $k$-means        & yes & spherical \\
    Ward linkage     & no  & compact \\
    DBSCAN           & no  & arbitrary \\
    \bottomrule
  \end{tabular}}
\end{frame}

\begin{frame}{Big takeaway}
  \begin{takeaway}
    The objective is the only teacher.\\
    Choose it before you choose the algorithm.
  \end{takeaway}
\end{frame}

\begin{frame}[fragile]{Code}
  Running $k$-means in scikit-learn:
  \vskip1ex
  \begin{tdcode}
\begin{Verbatim}
from sklearn.cluster import KMeans
km = KMeans(n_clusters=3, n_init=10)
labels = km.fit_predict(X)
print(km.inertia_)
\end{Verbatim}
  \end{tdcode}
\end{frame}

% ======================================================================
\section{Outlook}
% ======================================================================

\begin{frame}{Closing}
  Everything on this slide is serif Palatino, near-black canvas, DHBW red
  for the one thing that matters. The footer fills as we go.
  \begin{itemize}
    \item Margin notes are on-demand, not a standing column.
    \item Figures are recolored at the source, not inverted.
  \end{itemize}
\end{frame}

\end{document}
